% Original synthetic teaching sample for TeXada; AGPL-3.0-only.
% No copied textbook content or real personal data.
\documentclass[12pt]{article}
\usepackage{amsmath,amssymb}
\usepackage{listings}
\usepackage[margin=2.2cm]{geometry}
\setlength{\parindent}{0pt}
\setlength{\parskip}{0.7em}
% Teaching focus: TeX context decides what counts as math source.
% A dollar pair inside a comment, like $1+1=3$, never becomes a candidate.
% This commented mirror identity $a^2-b^2=(a-b)(a+b)$ is never re-checked.
\title{Context Notes: Comments, Verbatim and Display Math}
\author{LinguistsWantTech --- synthetic teaching sample}
\date{}
\begin{document}
\maketitle
\section*{1. Audited inline formulas}
The identity below is real math source:
$a^2-b^2=(a-b)(a+b)$.
A share line stays scannable across an escaped percent: the fraction
$s=\frac{x}{x+y}$ reaches 30\% of the sum $x+y$.
\section*{2. Deliberate syntax gap in real text}
One fraction below deliberately misses its closing brace:
$\displaystyle \frac{3x}{2x+1$.
\section*{3. Verbatim and listings bodies are literal}
\begin{verbatim}
Raw snippet: $1+1=3$ and \left( stay literal here.
\end{verbatim} Scanning resumes on this line: $a_{n+1}=q\,a_n$.
\begin{lstlisting}[language=Python]
rate = "$1+1=3$"  # dollar pairs stay code, not math
\end{lstlisting}
\section*{4. Same-line display math stays outside inline candidates}
$$\displaystyle g(t)=\left(\frac{2t}{1+t}$$
The sized parenthesis above is deliberately unclosed. Display math and
verbatim defects therefore need manual review, exactly like the align and
label defects recorded in the earlier boundary samples.
\end{document}
